\documentclass{amsart}
% For dealing with references we use the comment environment
\usepackage{verbatim}
\newenvironment{reference}{\comment}{\endcomment}
%\newenvironment{reference}{}{}
\newenvironment{slogan}{\comment}{\endcomment}
\newenvironment{history}{\comment}{\endcomment}

% For commutative diagrams we use Xy-pic
\usepackage[all]{xy}

% We use 2cell for 2-commutative diagrams.
\xyoption{2cell}
\UseAllTwocells

% We use multicol for the list of chapters between chapters
\usepackage{multicol}

% This is generally recommended for better output
\usepackage{lmodern}
\usepackage[T1]{fontenc}

% For cross-file-references

% Package for hypertext links:
\usepackage{hyperref}

% For any local file, say "hello.tex" you want to link to please

% Theorem environments.
%
\theoremstyle{plain}
\newtheorem{theorem}[subsection]{Theorem}
\newtheorem{proposition}[subsection]{Proposition}
\newtheorem{lemma}[subsection]{Lemma}

\theoremstyle{definition}
\newtheorem{definition}[subsection]{Definition}
\newtheorem{example}[subsection]{Example}
\newtheorem{exercise}[subsection]{Exercise}
\newtheorem{situation}[subsection]{Situation}

\theoremstyle{remark}
\newtheorem{remark}[subsection]{Remark}
\newtheorem{remarks}[subsection]{Remarks}

\numberwithin{equation}{subsection}

% Macros
%
\def\lim{\mathop{\mathrm{lim}}\nolimits}
\def\colim{\mathop{\mathrm{colim}}\nolimits}
\def\Spec{\mathop{\mathrm{Spec}}}
\def\Hom{\mathop{\mathrm{Hom}}\nolimits}
\def\Ext{\mathop{\mathrm{Ext}}\nolimits}
\def\SheafHom{\mathop{\mathcal{H}\!\mathit{om}}\nolimits}
\def\SheafExt{\mathop{\mathcal{E}\!\mathit{xt}}\nolimits}
\def\Sch{\mathit{Sch}}
\def\Mor{\mathop{\mathrm{Mor}}\nolimits}
\def\Ob{\mathop{\mathrm{Ob}}\nolimits}
\def\Sh{\mathop{\mathit{Sh}}\nolimits}
\def\NL{\mathop{N\!L}\nolimits}
\def\CH{\mathop{\mathrm{CH}}\nolimits}
\def\proetale{{pro\text{-}\acute{e}tale}}
\def\etale{{\acute{e}tale}}
\def\QCoh{\mathit{QCoh}}
\def\Ker{\mathop{\mathrm{Ker}}}
\def\Im{\mathop{\mathrm{Im}}}
\def\Coker{\mathop{\mathrm{Coker}}}
\def\Coim{\mathop{\mathrm{Coim}}}

% Boxtimes
%
\DeclareMathSymbol{\boxtimes}{\mathbin}{AMSa}{"02}

%
% Macros for moduli stacks/spaces
%
\def\QCohstack{\mathcal{QC}\!\mathit{oh}}
\def\Cohstack{\mathcal{C}\!\mathit{oh}}
\def\Spacesstack{\mathcal{S}\!\mathit{paces}}
\def\Quotfunctor{\mathrm{Quot}}
\def\Hilbfunctor{\mathrm{Hilb}}
\def\Curvesstack{\mathcal{C}\!\mathit{urves}}
\def\Polarizedstack{\mathcal{P}\!\mathit{olarized}}
\def\Complexesstack{\mathcal{C}\!\mathit{omplexes}}
% \Pic is the operator that assigns to X its picard group, usage \Pic(X)
% \Picardstack_{X/B} denotes the Picard stack of X over B
% \Picardfunctor_{X/B} denotes the Picard functor of X over B
\def\Pic{\mathop{\mathrm{Pic}}\nolimits}
\def\Picardstack{\mathcal{P}\!\mathit{ic}}
\def\Picardfunctor{\mathrm{Pic}}
\def\Deformationcategory{\mathcal{D}\!\mathit{ef}}

\setlength{\emergencystretch}{3em}
\begin{document}
\title[Stacks Project, Tag 07HD]{330 --- Divided power envelopes commute with Tor-controlled flat base change}
\maketitle
\noindent Copyright (C) 2005--2025 Johan de Jong. Stacks Project authors. Source: \href{https://stacks.math.columbia.edu/tag/07HD}{Tag 07HD}.

\noindent Editorial difficulty note (not author text). Difficulty H2: The proof must extend divided powers through a possibly nonflat algebra base change. A polynomial presentation and the Tor hypothesis identify the relation ideal intersection, while flatness modulo the original divided power ideal proves an auxiliary injection that allows divided powers to descend through the quotient..

\noindent Editorial provenance note (not author text). History: excerpt from revision \texttt{a0444\allowbreak 6e57e\allowbreak c1fbc\allowbreak 252a8\allowbreak 71afc\allowbreak ec775\allowbreak 2fb28\allowbreak 07b14}, crystalline.tex, lines 285--362. Reference presentation was adapted; original-author context and core support blocks were retained or added. The mathematical wording is unchanged. Licensed under GNU FDL 1.2 or later; no invariant sections or cover texts. See ../licenses/COPYING. Context blocks reproduce author definitions and supporting results; they are not additional counted problems.

\begin{lemma}
\label{lemma-flat-base-change-divided-power-envelope}
Let $(A, I, \gamma)$ be a divided power ring.
Let $B \to B'$ be a homomorphism of $A$-algebras.
Assume that
\begin{enumerate}
\item $B/IB \to B'/IB'$ is flat, and
\item $\text{Tor}_1^B(B', B/IB) = 0$.
\end{enumerate}
Then for any ideal $IB \subset J \subset B$ the canonical map
$$
D_B(J) \otimes_B B' \longrightarrow D_{B'}(JB')
$$
is an isomorphism.
\end{lemma}

\begin{proof}
Set $D = D_B(J)$ and denote $\bar J \subset D$ its divided power ideal
with divided power structure $\bar\gamma$. The universal property of
$D$ produces a $B$-algebra map $D \to D_{B'}(JB')$, whence a map as in
the lemma. It suffices to show that
the divided powers $\bar\gamma$ extend to $D \otimes_B B'$ since then
the universal property of $D_{B'}(JB')$ will produce a map
$D_{B'}(JB') \to D \otimes_B B'$ inverse to the one in the lemma.

\medskip\noindent
Choose a surjection $P \to B'$ where $P$ is a polynomial algebra over $B$.
In particular $B \to P$ is flat, hence $D \to D \otimes_B P$ is flat by
Algebra, Lemma \href{https://stacks.math.columbia.edu/tag/00HI}{lemma flat base change (Tag 00HI)}.
Then $\bar\gamma$ extends to $D \otimes_B P$ by
Divided Power Algebra, Lemma \href{https://stacks.math.columbia.edu/tag/07H1}{lemma gamma extends (Tag 07H1)}; we
will denote this extension
$\bar\gamma$ also. Set $\mathfrak a = \Ker(P \to B')$ so that
we have the short exact sequence
$$
0 \to \mathfrak a \to P \to B' \to 0
$$
Thus $\text{Tor}_1^B(B', B/IB) = 0$ implies that
$\mathfrak a \cap IP = I\mathfrak a$.
Now we have the following commutative diagram
$$
\xymatrix{
B/J \otimes_B \mathfrak a \ar[r]_\beta &
B/J \otimes_B P \ar[r] &
B/J \otimes_B B' \\
D \otimes_B \mathfrak a \ar[r]^\alpha \ar[u] &
D \otimes_B P \ar[r] \ar[u] &
D \otimes_B B' \ar[u] \\
\bar J \otimes_B \mathfrak a \ar[r] \ar[u] &
\bar J \otimes_B P \ar[r] \ar[u] &
\bar J \otimes_B B' \ar[u]
}
$$
This diagram is exact even with $0$'s added at the top and the right.
We have to show the divided powers on the ideal
$\bar J \otimes_B P$ preserve the ideal
$\Im(\alpha) \cap \bar J \otimes_B P$, see
Divided Power Algebra, Lemma \href{https://stacks.math.columbia.edu/tag/07H2}{lemma kernel (Tag 07H2)}.
Consider the exact sequence
$$
0 \to \mathfrak a/I\mathfrak a \to P/IP \to B'/IB' \to 0
$$
(which uses that $\mathfrak a \cap IP = I\mathfrak a$ as seen above).
As $B'/IB'$ is flat over $B/IB$ this sequence remains exact after
applying $B/J \otimes_{B/IB} -$, see
Algebra, Lemma \href{https://stacks.math.columbia.edu/tag/00HL}{lemma flat tor zero (Tag 00HL)}. Hence
$$
\Ker(B/J \otimes_{B/IB} \mathfrak a/I\mathfrak a \to
B/J \otimes_{B/IB} P/IP) =
\Ker(\mathfrak a/J\mathfrak a \to P/JP)
$$
is zero. Thus $\beta$ is injective. It follows that
$\Im(\alpha) \cap \bar J \otimes_B P$ is the
image of $\bar J \otimes \mathfrak a$. Now if
$f \in \bar J$ and $a \in \mathfrak a$, then
$\bar\gamma_n(f \otimes a) = \bar\gamma_n(f) \otimes a^n$
hence the result is clear.
\end{proof}
\end{document}
